\chapter{스핑크스와 주피터}

\term{스핑크스}\annotate{Sphinx}와 \term{주피터}\annotate{Jupyter}는 \term{파이썬}\annotate{Python}에 포함된 문서화 도구이다.

\section{스핑크스}

스핑크스가 생성한 텍 파일로부터 얻은 PDF에 한글 텍스트가 식자되지 않은 이유는 스핑크스가 한국어 텍을 위한 설정을 갖고 있지 않기 때문이다.

스핑크스는 텍 엔진에 따라 \macro{babel} 또는 \macro{polyglossia}를 사용하는 반면, "그림"이나 "경고" 같은 이름들에 대해 \macro{polyglossia}에 딸린 언어 정의 파일들이 아니라 스핑크스가 갖고 있는 로캘 파일들을 사용한다. 
그래서 \macro{kotex}이 사용되도록 조작하는 것만으로 문제를 완전히 해결할 수 없다.

\macro{conf.py}가 다음과 같이 작성되어야 한다.

\begin{codewrite}
from docutils.writers.latex2e import Babel
Babel.language_codes["ko"]="korean"
language='ko'
latex_engine = 'xelatex'
latex_documents = [(master_doc, 'foo.tex', project, author, 'hzguide')]
latex_elements = {
    'babel': ' ',
    'preamble': r'''
        \usepackage{kotex}
        \usepackage{…}
        \setmainfont{…}
        \setmainhangulfont{…}
        '''
}
\end{codewrite}
\pythonread

스핑크스가 생성한 텍 문서에 사용된 명령들과 환경들은 스핑크스가 제공하는 \macro{sphinx.sty} 패키지에 정의되어 있다.
대부분은 재정의하여 스타일을 변경할 수 있지만, 절 스타일은 \macro{titlesec} 패키지를 이용하여 정의되어 있기 때문에 HzGuide가 제어할 수 없다.
해법은 \macro{sphinx} 패키지를 불러들이기 전에 \macro{titlesec}를 무력화하는 것이다.

\begin{code}
\EmulatedPackageWithOptions{nobottomtitles*}{titlesec}[2019/10/16]
\DeclareDocumentCommand\titleformat{ smommmmo }{}
\EmulatedPackage{sphinxmessages}
\end{code}

\section{주피터}

주피터 파일로부터 레이텍을 거쳐 만들어진 PDF에서도, 스핑크스에서와 마찬가지로, 한글 텍스트가 식자되지 않는 문제가 있다.
손쉬운 해법은 다음 경로에 있는 템플릿 파일을 참고하여 새 템플릿을 만드는 것이다.

\begin{codewrite}
C:/…/Lib/site-packages/nbconvert/templates/latex/document_contents.tplx
\end{codewrite}
\powershellread

새 템플릿은 다음과 같은 코드를 포함해야 한다.

\begin{codewrite}
((* block docclass *))
    \documentclass{hzguide}
    \usepackage{…}
((* endblock docclass *))
\end{codewrite}
\coderead[language=js+jinja]

\begin{codewrite}
C:\>jupyter nbconvert --to latex --template FOO.tplx foo.ipynb
\end{codewrite}
\powershellread